\documentclass[12pt,a4paper]{report}

\usepackage[utf8]{inputenc}
\usepackage[T2A]{fontenc}
\usepackage[serbianc]{babel}
\usepackage{amsmath, amssymb}
\usepackage{graphicx}
\usepackage{caption}
\usepackage{geometry}
\geometry{margin=2.5cm}

\begin{document}

\subsection{Тахографи}

Тахографи тзв. региструјући инструменти са писаљкама којима се бележе: брзина возила, пређени пут, време вожње и заустављања и остале важније информације. Уграђују се у камионе и аутобусе заједно са брзиномером и километар-сатом.

\textbf{Папир за регистровање} је кружног облика са исцртаним линијама дијаграма по коме писаљка инструмента оставља траг. С обзиром на то да се подаци бележе у зависности од времена, у овај инструмент уграђен је и сатни механизам који помера папир за регистровање, обично цео круг за 24 сата (види Слика 3.82).

По спољашњем и унутрашњем ободу дијаграма налази се сатна скала са поделом чији је размак од 5 минута. Између ове две скале је скала путне брзине у облику концентричних кругова са поделом у размаку од $20\text{ km/h}$. Вредност брзине се записује тако што се писаљка радијално помера (ка спољнем ободу) у зависности од брзине, а сатни механизам помера папир за регистровање у круг, сходно временској скали.

У унутрашњем делу папира уписује се пређени пут у облику троугластих осцилација, при чему сваком „троуглу“ одговара пут од $10\text{ km}$. Праве линије између њих означавају време заустављања возила. Са спољне стране овог круга дебелом линијом обележава се време вожње а

% --- Страница 194 ---
танком линијом време стајања возила. Унутар круга ручно се уписују основни подаци о возачу, полазно место и одредиште пута, датуми и пређени километри.

\textbf{Индукциони инструмент.} – Код механичких тахографа, скоро искључиво користе се индукциони инструменти, као и код механичких брзиномера (слика 3.69). Разлика је само у томе што се за тахограф (као региструјући инструмент са писаљком) користи много јачи стални магнет који ствара већи обртни момент. Мерно звоно, помоћу малог зупчаника и зупчасте летве, помера иглу са сафирним врхом који по импрегнираном папиру уписује (гребањем) податке о кретању возила.

Линија пређеног пута у облику шпицастих осцилација уписује се помоћу посебне писаљке коју помера мала осовина са једним брегом.

Осовина се помера преко механизма бројчаника километар-сата, тако да се обрне за један цео круг за пређених $10\text{ km}$.

% \begin{figure}[htbp]
%     \centering
%     \includegraphics[width=0.5\textwidth]{slika_3_82.png}
%     \caption{Облик дијаграма на папиру за регистровање у тахографу}
%     \label{fig:3.82}
% \end{figure}

\subsubsection{Електрични (електронски) тахографи}

Код механичких тахографа радни век савитљиве погонске осовине (сајле) је кратак због великог статичког и динамичког оптерећења и релативно велике дужине. Откуда се и овде примењује решење са електричним покретањем.

Код електричних тахографа, наместо индукционог мерног система, управља се електрични. Слично електричним брзиномерима и овде се
% --- Страница 195 ---
механички обртаји најпре претварају у електрични сигнал, који се у мерном инструменту поново претвара у механички момент који покреће систем за кретање инструмента и писаљке.

\textbf{Принцип рада.} – Код електричног тахографа, као и код брзиномера, неопходна су два одвојена претварача сигнала за брзину и пређени пут. То је неопходно јер електрични сигнал брзиномера има велику учестаност, а за покретање електромеханичког бројчаника и регистратора, неопходан је сигнал мале учестаности.

Брзиномеру импулсе даје мали електрични генератор, чији ротор од сталног магнета има 36 полова. Кавезасти намотај статора, такође са 36 полова, даје наизменичну струју са 18 периода по обртају ротора. Најнижа учестаност за почетак рада брзиномера, је око $15\text{ Hz}$, која се постиже већ при брзини возила од око $3\text{ km/h}$.

Давач импулса за пређени пут је најчешће тзв. Reed-контактни претварач који за сваки обртај даје један импулс, тако да су поуздани и при малим брзинама. Овај давач се електриком напаја преко акумулатора. Сигнали давача доводе се до малог степ-мотора који у инструменту покреће бројач пређених километара и механизам тахографа са писаљком.

На слици 3.83 шематски је приказан принцип рада електричног тахографа са брзиномером.

% \begin{figure}[htbp]
%     \centering
%     \includegraphics[width=0.7\textwidth]{slika_3_83.png}
%     \caption{Принцип рада електричног тахографа са брзиномером (Kienzle): 1. давач, 2. обликовање сигнала, 3. конструкција електронског брзиномера, 4. диференцијатор сигнала, 5. појачивач, 6. серво-мотор, 7. повратна спрега, 8. сензор угла померања, 9. управљачко коло, 10. зупчаник, 11. зупчаста летва, 12. писаљка, 13. казаљка брзиномера, 14. скала}
%     \label{fig:3.83}
% \end{figure}

Електрични сигнал давача се мора обликовати и по величини напона и по учестаности. Електрична веза, слична електронском брзиномеру, помоћу Schmitt-trigger-а најпре обликује четвртасте импулсе, који се затим, моностабилним мултивибратором, униформишу у импулсе константне ширине и амплитуде. Просечна збирна вредност тих импулса представља електрични импулс, сразмеран брзини, који се за покретање писача и казаљке инструмента мора поново претворити у механички сигнал.

Казаљку и писач брзиномера покреће мали једносмерни серво-мотор, који се напаја преко појачивача и управљачког кола. Тиме се обезбеђује да угао скретања мотора буде сразмеран напону. Стварно скретање региструје мали потенциометар повезан са осовином, чији се сигнал враћа на улазни диференцијатор. Евентуално одступање између
% --- Страница 196 ---
задатог и стварног скретања, у облику електричног сигнала, преко појачивача, враћа се у управљачко коло.

На тај начин кретни систем ће се умирити само када је то одступање једнако нули. Тачност скретања још више се повећава преко повратне спреге (7), којом се обезбеђује и боља статичка и динамичка регулација система.

Обртање папира за регистровање са дијаграмом врши сатни механизам са кварцним осцилатором ради веће тачности.


\subsection{Мерачи притиска – манометри}

За разлику од брзиномера и километар-сата, који дају податке о кретању возила, мерачи притиска спадају у групу инструмената који дају податке о погонској ситуацији возила. Код мотора са унутрашњим сагоревањем (SUS-мотори) ови мерачи најчешће прате притисак уља за подмазивање или потпритисак у усисној цеви, а код кочионог система притисак главног цилиндра или кочионих цилиндара на точковима.

Као и сви остали инструменти и они се састоје од два дела – давача и показног инструмента. Давачи су у ствари претварачи сигнала, који се увек користе за мерење неелектричних величина електричним путем. Уграђују се непосредно на место где се мери притисак и у директном су додиру са мереном величином.

За мерење притиска најчешће се користе давачи са биметалном траком или потенциометром (слика 3.84).

% \begin{figure}[htbp]
%     \centering
%     \includegraphics[width=0.6\textwidth]{slika_3_84.png}
%     \caption{Шематски приказ давача за мерење притиска: а) са биметалом и б) са потенциометром}
%     \label{fig:3.84}
% \end{figure}

Као што се са слике види, доњи део давача је од масивног метала, који се уврће у простор под притиском, тако да тај притисак делује на еластичну мембрану. У горњем делу је електрични део давача, који је најчешће од биметалне траке са електричним загревањем, или потенциометра са променљивим отпорником.

\textbf{Биметални давач} се састоји од горње биметалне траке са загревним влакном и доње траке у споју са еластичном мембраном. Помоћу својих контаката траке су међусобно спојене. Када се на давач, преко показног инструмента, прикључи електрични напон, једна мала струја протиче кроз загревно влакно, контакте и доње траке према маси. Услед загревања горња биметална трака се извије нагоре, раздвоји контакте и прекине струју. По прекиду струје биметална трака се охлади и поново споји контакте. На тај начин кроз
% --- Страница 197 ---
давач и инструмент тече испрекидана струја, с тим што ритам прекидања зависи од притиска мембране на доњу траку и њен контакт. Када је притисак једнак нули, контакти се лако и брзо отварају па су струјни импулси краткотрајни и ретки (види слику 3.85 – положај за притисак $0\text{ bar}$).

Ако притисак порасте, на пример на $2\text{ бара}$, мембрана потисне навише доњи контакт за ($h_1$). У том случају, за раздвајање контаката је потребно веће извијање биметала, што се постиже већим загревањем са чешћим струјним импулсима (слика 3.85 – положај за притисак од $2\text{ бара}$).

% \begin{figure}[htbp]
%     \centering
%     \includegraphics[width=0.75\textwidth]{slika_3_85.png}
%     \caption{Принцип рада биметалног мерача притиска са дијаграмом мерене струје у зависности од притиска: $h_1$ и $h_2$ – висина подизања контаката, $I_0$ – величина пулсирајуће струје, $I_{\text{ef}}$ – ефективна вредност пулсирајуће струје, $T_z$ – време затварања контаката, $T_o$ – време отварања контаката}
%     \label{fig:3.85}
% \end{figure}

При још већем притиску, мембрана ће још више потиснути контакте навише, па је за њихово раздвајање потребно још веће загревање и чешћи струјни импулси (слика 3.85 – положај за притисак од $5\text{ бара}$).

\textbf{Биметални инструмент.} – На доњем делу слике (3.85) симболично је приказано и скретање биметалног показног инструмента, у зависности од учесталости струјних импулса и њене ефективне вредности. Принцип рада и конструкција овог инструмента шематски су приказани на слици 3.86.

Основни део инструмента је биметал (4) у облику слова „U“ на коме је загревно влакно (3) од танке отпорне жице. Један крак биметала је спојен за лимени држач са зупцима за подешавање (11), а други, зглобно, за казаљку (9). Лимени држач (11) заједно са биметалом, помоћу назубљеног лука, закреће се око осовине (1).

Други лимени подешивач (6) преко еластичне траке (7) такође подупире зглавкасти спој казаљке у лежишту (8). Његов положај по-%
% --- Страница 198 ---
% \begin{figure}[htbp]
%     \centering
%     \includegraphics[width=0.6\textwidth]{slika_3_86.png}
%     \caption{Принцип рада и конструкција биметалног инструмента: 1. осовина лименог држача, 2. биметал у облику слова U, 3. влакна која се загревају, 4. назубљени лук подешивача, 5. прикључни контакти за довод струје, 6. лимени подешивач, 7. еластична трака, 8. лежиште казаљке, 9. казаљка, 10. неузубљени лук, 11. лимени подешивач држача}
%     \label{fig:3.86}
% \end{figure}
дешава се помоћу назубљеног лука (4). Први крак биметала, који је у споју са држачем (11) врши температурну компензацију, а други радни крак, на коме је загревно влакно, помера казаљку инструмента. Електрично напајање загревног влакна (3) изведено је преко прикључка (5). Један прикључак се везује за биметални давач а други за позитивни пол акумулатора (преко контакт-кључа). Скретање казаљке зависи од величине ефективне вредности струјних импулса, како је то симболично приказано на слици 3.85.

\textbf{Температурна компензација.} – Померање биметалне траке зависи и од температуре околине, што неповољно утиче на тачност мерења. Да би се умањио тај утицај биметалне траке се израђују са два крака у облику слова U, од којих је један радни (у споју са казаљком) а други компензациони (у споју са постољем) (слика 3.87).

% \begin{figure}[htbp]
%     \centering
%     \includegraphics[width=0.6\textwidth]{slika_3_87.png}
%     \caption{Померање биметалне траке у облику слова U под утицајем температуре околине}
%     \label{fig:3.87}
% \end{figure}

У почетном положају (I) на општој температури, оба крака биметалне траке су паралелна. Са порастом температуре околине, компензациони крак се криви удесно, у односу на доњу фиксну тачку, померајући и радни крак, који се, у овом примеру, уопште не ради и без струјног загревања. Међутим, повећана температура околине утиче и на радни крак, па се и он криви удесно, али у односу на горњу спојну тачку (положај III). На тај начин, без обзира на померање спојне тачке ($B_1$), тачка (A), која је у споју са казаљком инструмента, остала је релативно непомична у односу на положај (I).

\textbf{Напонска компензација.} – На тачност показивања биметалних инструмената и давача битно утиче и стабилност напона напајања. Из тих разлога у ове инструменте обавезно се уграђују мали биметални или електронски стабилизатори напона.
% --- Страница 199 ---
\textbf{Показни инструмент.} – Код давача са потенциометром, као показни инструмент за мерење притиска највише се користе инструменти са укрштеним навојима и обртним магнетом, тзв. логометри. Карактеристика ових инструмената је да немају отпорни момент који ствара спирална опруга, већ сваки навој производи свој обртни момент који су међусобно супротних смерова. Систем кретања се зауставља у положају када се ти моменти изједначе.

На возилама ови инструменти имају три непомична намотаја који су међусобно под правим углом. Унутар навоја је лако обртан, попречно намагнетисан стални магнет са казаљком. Магнет се обрће под дејством магнетних поља намотаја и заузима положај резултујућег магнетног поља. На слици 3.88 дата је шема веза намотаја са регулационим отпорницима и уцртаним смеровима струје и магнетних поља. Потенциометарски давач приказан је у облику променљивог отпорника ($R_a$).

% \begin{figure}[htbp]
%     \centering
%     \includegraphics[width=0.75\textwidth]{slika_3_88.png}
%     \caption{Шема везе инструмента са укрштеним навојима и обртним магнетом за мерење притиска са потенциометарским давачем}
%     \label{fig:3.88}
% \end{figure}

У почетном положају, када је мерени притисак једнак нули, отпор потенциометра давача је максималан и кроз трећи намотај тече незнатно мала струја. У том случају струја из акумулатора протиче кроз први и други намотај и обртни магнет се постави у положај резултујућег магнетног поља тих намотаја ($H_{1,2}$), па је то уједно и нулти положај инструмента (слика 3.88). Када се мерени притисак повећа, отпор потенциометра давача се смањи а струја у трећем намотају и магнетно поље ($H_3$) повећају. Тада, на обртни магнет, поред резултујућег магнетног поља ($H_{1,2}$), које је константно, делује и магнетно поље $H_3$. Магнет са казаљком инструмента ће се обрнути у правцу резултујућег поља ($H_r$). Ово скретање је сразмерно мереном притиску који помоћу еластичне мембране мења отпор давача ($R_a$).
% --- Страница 200 ---
Отпорник ($R_3$) служи за температурну компензацију промене отпора намотаја, а отпорником ($R_1$) регулише се јачина струје у намотајима. Конструкција оваквог инструмента приказана је на слици 3.89.

% \begin{figure}[htbp]
%     \centering
%     \includegraphics[width=0.7\textwidth]{slika_3_89.png}
%     \caption{Конструкција инструмента чија је шема везе дата на слици (3.88): 1. постоље, 2. осовина обртног магнета, 3. завртањ лежишта, 4. велики калем са 13 намотаја, 5. велики калем са 13 намотаја, 6. заштитни омотач, 7. магнет који казаљку држи у нултом положају, 8. мали бочни калем, 9. крстасти држач, 10. завртањ, 11. матица лежишта, 12. матица, 13. казаљка са контратеговима}
%     \label{fig:3.89}
% \end{figure}

\subsection{Термометри за мерење температуре расхладне течности}

За сигуран и економичан рад мотора температура расхладне течности мора се одржавати у одређеним границама, што зависи од конструкције и типа мотора. Због тога температура расхладне течности аутоматски се регулише преко расхладног система и истовремено се помоћу мерног инструмента возач информише о тој температури.

Температура неког медија може се мерити механички, на основу ширења неке течности или електрично, помоћу разних претварача, који-%
% --- Страница 201 ---
ма се промена температуре претвара у електрични сигнал. У даљем тексту обрађени су само електрични термометри.

\textbf{Принцип рада.} – Принцип рада електричних термометара сличан је раду описаних мерача притиска. Код биметалних инструмената користе се такође биметални давачи, а код инструмената са укрштеним навојем, давачи су најчешће отпорници чија се отпорност мења са температуром.

Из основа електротехнике је познато да код већине метала електрична отпорност расте са температуром. Код одређених полупроводника то је обрнуто, јер њихова отпорност опада са порастом температуре.

Као температурни давачи, највише се користе полупроводници од керамичких једињења са оксидима тешких метала, познатији под називом „термистори“.

\textbf{Термометри са укрштеним навојем.} – Најједноставнија варијанта ових инструмената има два укрштена навоја под правим углом (слика 3.90).

% \begin{figure}[htbp]
%     \centering
%     \includegraphics[width=0.75\textwidth]{slika_3_90.png}
%     \caption{Шема везе инструмента са два укрштена навоја и термисторским давачем}
%     \label{fig:3.90}
% \end{figure}

На слици 3.90 а приказана је електрична шема, према стварном распореду намотаја, предпротора, компензационог отпорника и отпорника давача, а на слици 3.90 б ради лакшег разумевања, дата је иста шема у тзв. развијеном облику.

Електрична отпорност давача ($R_a$), са негативним температурним коефицијентом (термистор), везан је паралелно са другим намотајем инструмента ($L_2$). Отпорност термистора осетно опада са порастом температуре (види табелу).

\begin{center}
\begin{tabular}{lccc}
\hline
Температура ($^\circ\text{C}$) & 40 & 90 & 120 \\
Отпорност ($\Omega$) & 137 & 23,5 & 10,2 \\
\hline
\end{tabular}
\end{center}

Када је температура течности мала, електрични отпор давача ($R_a$) је велик, па је струја у паралелно везаном намотају ($L_2$) скоро иста као и у редно везаном намотају ($L_1$). У том случају обртни део инструмента од
% --- Страница 202 ---
меког гвожђа, са казаљком, се поставља у правцу резултујућег магнетног поља ($H_r$) где је поље ($H_{2k}$) другог калема велико а поље ($H_{1k}$) првог калема мало (слика 3.90б).

Са порастом температуре течности, опада отпорност давача ($R_a$), а са тим и пад напона на паралелном делу везе (слика 3.90б). Услед тога, струја у првом калему ($L_1$) порасте а са њом и магнетно поље ($H_1$), док струја у другом калему ($L_2$) опадне а са њом се смањи и његово магнетно поље ($H_2$). У том случају, обртно меко гвожђе, заједно са казаљком, помериће се улево, у правцу новог резултујућег поља ($H_{rv}$). Највеће скретање инструмента и ширина скале су релативно мали, што и јесте највећа мана ових инструмената.

На слици 3.91 дата је шема везе инструмента са три укрштена навоја који је по конструкцији сличан инструменту за мерење притиска (слика 3.88).

% \begin{figure}[htbp]
%     \centering
%     \includegraphics[width=0.5\textwidth]{slika_3_91.png}
%     \caption{Шема везе инструмента са три укрштена навоја и термисторским давачем}
%     \label{fig:3.91}
% \end{figure}

Овде је на ред са давачем везан трећи намотај (III). Отпорником ($R_e$) се подешава напон напајања а отпорником ($R_s$) скретање инструмента, у зависности од опсега промене отпора давача.

У мерној грани, коју сачињавају први и други намотај постављени под правим углом, протиче константна струја која не зависи од промене отпора давача. На тај начин, резултујуће магнетно поље зависи само од струје у трећем намотају, тј. од промене отпора давача. Отпорност давача мења се у границама мерног подручја инструмента од $40$ до $120^\circ\text{C}$ а ширина скале је $90^\circ$.

\textbf{Конструкција биметалног давача.} – На слици 3.92 приказана је конструкција једног биметалног давача за мерење температуре.

% \begin{figure}[htbp]
%     \centering
%     \includegraphics[width=0.6\textwidth]{slika_3_92.png}
%     \caption{Конструкција биметалног давача за мерење температуре: 1. кућиште, 2. загревно влакно, 3. биметална трака, 4. помични контакт, 5. непомични контакт, 6. месингана чаура, 7. изолациони држач, 8. довод струје, 9. изолатор, 10. електрични прикључак}
%     \label{fig:3.92}
% \end{figure}

% --- Страница 203 ---
Механизам давача смештен је у месингану чауру (6), која је углављена у кућиште (1) у облику шестоугаоног завртња којима је давач учвршћен на месту уградње. Сигнал се претвара помоћу биметалне траке (3), која је са једне стране причвршћена за изолациони држач (7) а на другој страни је постављен помични контакт (4). Биметална трака се загрева помоћу влакна (2), чији је један крај спојен за сам биметал а други преко еластичног проводника (8), везан је за прикључак (10). Електрични довод је помоћу изолатора (9) изолован од металног кућишта (1). Крај загревног влакна, који је спојен за биметалну траку, преко непомичног контакта (5) спаја се са масом, тј. негативним полом акумулатора. Електрични прикључак давача (10) повезује се са биметалним инструментом чија је конструкција приказана на слици 3.86.

По укључењу контакт-кључа кроз мерни систем протиче струја и загревно влакно загрева биметалну траку. У почетном положају, кад је мерена течност хладна, биметал се спорије загрева, па се и контакти спорије раздвајају, што повећава ефективну вредност струјних импулса. Са порастом температуре течности (воде), преко месингане чауре, биметал се додатно загрева, па се и контакти брже отварају, те се смањује ефективна вредност струјних импулса.

На слици 3.93 шематски је приказан рад биметалног давача са биметалним инструментом. Као што се са слике види, инструмент ради у обрнутом смеру, те при малим температурама течности – скретање је велико а при великим температурама течности – скретање је мало. То је донекле и повољније јер је током вожње температура расхладне течности велика. У искљученом положају казаљка инструмента стоји на $100^\circ$, тј. приближно на десном крају скале.

% \begin{figure}[htbp]
%     \centering
%     \includegraphics[width=0.6\textwidth]{slika_3_93.png}
%     \caption{Шема везе биметалног давача са биметалним инструментом за мерење температуре расхладне течности}
%     \label{fig:3.93}
% \end{figure}

\textbf{Термисторски давач.} – На слици 3.94 дата је конструкција термисторског давача, који се најчешће везује са инструментом са укрштеним
% --- Страница 204 ---
навојима и обртним магнетом, или инструментима са електронским напајањем.

% \begin{figure}[htbp]
%     \centering
%     \includegraphics[width=0.5\textwidth]{slika_3_94.png}
%     \caption{Попречни пресек термисторског давача за мерење температуре: 1. месингано кућиште, 2. опружни проводник струје, 3. изолациона цев, 4. тањираста опруга, 5. плоча термистора, 6. изолатор}
%     \label{fig:3.94}
% \end{figure}

Термисторску плочицу (5), преко тањирасте опруге (4), опружни проводник струје (2) притиска на дно месинганог кућишта (1) која је у споју са масом. Остале унутрашње делове танка изолациона цев одваја од кућишта. Електрични отпор термометарског давача мења се у зависности од температуре од $1\,600$ до $60\text{ ohm-а}$.

\textbf{Давач са отпорном жицом.} – На слици 3.95 приказана је конструкција температурног давача са отпорном металном жицом, која је намотана на телу (2) и смештена у месингану чауру (5). Прикључена струја протиче преко опруге (3) и променљивог отпорника ка месинганом кућишту и маси.

% \begin{figure}[htbp]
%     \centering
%     \includegraphics[width=0.6\textwidth]{slika_3_95.png}
%     \caption{Термички давач са отпорном металном жицом за камионе: 1. изолација, 2. метални отпорник, 3. проводна опруга, 4. заштита, 5. месингано кућиште, 6. заштитна капица}
%     \label{fig:3.95}
% \end{figure}

Електрични отпор давача при струји од $40\text{ mA}$ на температури од $80^\circ\text{C}$ износи $148{,}2 \pm 8\text{ ohm-а}$.

\textbf{Електронски термометар.} – Осим ових класичних инструмената са укрштеним навојима, користе се и инструменти са електронским напајањем и термисторским давачем (слика 3.96).

% \begin{figure}[htbp]
%     \centering
%     \includegraphics[width=0.6\textwidth]{slika_3_96.png}
%     \caption{Шема везе електронског термометра}
%     \label{fig:3.96}
% \end{figure}

Када се укључи контакт-кључ (K), преко отпорника ($R_V$) потенциометар (P), термистора (Th) и отпора $R_2$ затвара се струјно коло. Услед те струје јавља се пад напона на потенциометру и база транзистора (T) постаје негативна у односу на емитер. Због тога ће потећи базна стру-%
% --- Страница 205 ---
ја која ће отворити транзистор и укључити инструмент у струјно коло емитера и колектора. Јачина те струје зависи од величине базне струје, а базна струја од пада напона на потенциометру. Пошто тај пад напона зависи од јачине струје у колу са термистором, јасно је да ће од отпора термистора директно утицати на скретање инструмента. Отпор термистора, како је већ објашњено, мења се у зависности од температуре расхладне течности. Из тога произлази да је јачина струје кроз инструмент, преко транзисторског појачавача, директно сразмерна температури расхладне течности.

\subsection{Показивачи нивоа горива}

Ови инструменти сигнализирају возачу о стању горива у резервоару. Њихова тачност показивања је релативно мала. Од њих се само захтева да показују тренутну количину горива у односу на укупну запремину резервоара. Зато је и подела на скали најчешће означена бројевима: $0, 1/4, 1/2, 3/4, 1$.

Код појединих возила, на самом инструменту, улази се мала сигнална лампица, која упозорава, да на најнижи достигнут ниво горива у резервоару, тзв. резерву. Са том количином возач тачно зна колико још километара може прећи.

\textbf{Биметални показивачи.} – Код великог броја возила још увек се користе биметални показивачи нивоа горива или инструменти са укрштеним навојима и обртним магнетом. Биметални инструмент се повезује са биметалним или потенциометарским давачем, а инструмент са укрштеним навојима само са потенциометарским давачем.

Принцип рада ових показивача сличан је биметалним инструментима за мерење притиска (слика 3.85) с тим што контакте давача овде потискује пловак у резервоару а не притисак уља у мотору. Како је сила потиска механизма са пловком сразмерна нивоу горива у резервоару и ефективна вредност електричних импулса биметалног давача биће сразмерна том нивоу. Највећа вредност тих импулса је при пуном резервоару, зато је и скретање показног инструмента тада највеће. Принципијелна шема везе оваквог показивача дата је на слици 3.97.

% \begin{figure}[htbp]
%     \centering
%     \includegraphics[width=0.6\textwidth]{slika_3_97.png}
%     \caption{Принципијелна шема везе биметалног показивача нивоа горива (VDO)}
%     \label{fig:3.97}
% \end{figure}

Она је иста као и код биметалних термометара, с тим што се исти стабилизатор може користити за више мерних система.

\textbf{Показивачи са укрштеним навојима.} – Код путничких возила, ниже и средње класе, доста су распрострањени показивачи нивоа горива, код којих је инструмент са укрштеним навојима, а давач потенциометар са клизачем који помера пловак. Принципијелна шема везе дата је на слици (3.98).

Отпорници $R_1$ и $R_2$ су променљиви делови истог потенциометра који је смештен у резервоар. Зато је њихова збирна вредност увек иста без
% --- Страница 206 ---
обзира на положај клизача, јер померањем клизача мења се само њихов однос.

% \begin{figure}[htbp]
%     \centering
%     \includegraphics[width=0.75\textwidth]{slika_3_98.png}
%     \caption{Принципијелна шема везе показивача нивоа горива са укрштеним навојима: а) принципијелна шема, б) развијена шема, с) дијаграм магнетног поља}
%     \label{fig:3.98}
% \end{figure}

Тако, на пример, први отпорник ($R_1$), представља отпорност од почетка потенциометра до средњег клизача, а други ($R_2$), отпорност од клизача до краја потенциометра. Значи, померањем клизача један отпорник се повећава а други сразмерно смањује.

Код овог инструмента, за разлику од термометра на слици 3.90, намотаји инструмента су међусобно паралелно везани, а са сваким од њих по један отпорник на ред. Ако је резервоар празан, тада пловак падне најниже, па је у том положају клизача, на пример, вредност отпорника ($R_1$) највећа, а отпорника ($R_2$) најмања. Услед тога ће струја у првом навоју ($L_1$) бити највећа, а у другом ($L_2$) најмања. Зато ће магнетно поље првог навоја ($H_{1p}$) бити знатно веће од поља другог навоја ($H_{2p}$). Помично меко гвожђе са казаљком инструмента поставиће се у правац резултујућег поља ($H_{rp}$), тј. на почетак скале (слика 3.98).

Код пуног резервоара све се симетрично мења. Пловак се подигне на највиши положај, вредност отпора ($R_2$) порасте а отпора $R_1$ смањи, па се сходно томе, струја у навоју ($L_1$) смањи а у ($L_2$) повећа. У том случају магнетно поље другог навоја ($H_{2p}$) је знатно веће а првог навоја ($H_{1p}$) мање, па се меко гвожђе са казаљком постави у правцу новог резултујућег поља ($H_{rp}$), тј. на крај скале.

Иако је угао између намотаја и њихових магнетних поља $90^\circ$ угао скретања инструмента ($\beta$), између два резултујућа поља је свега $45^\circ$, што представља прилично узану скалу. Међутим, с обзиром на то да се ради само о показивању нивоа горива (а не и о тачном мерењу неке количине) скала је сасвим довољна за исписивање горе наведених пет бројева ($0, 1/4, 1/2, 3/4$ и $1$).

Казаљка инструмента у искљученом положају враћа се на нулти положај помоћу силе земљине теже, зато је радни положај инструмента одређен.
% --- Страница 207 ---
На слици 3.99 је дата конструкција и шематски приказ принципа рада инструмената са укрштеним навојима и помичним меким гвожђем.

% \begin{figure}[htbp]
%     \centering
%     \includegraphics[width=0.75\textwidth]{slika_3_99.png}
%     \caption{Конструкција и шема везе инструмента са укрштеним навојима и помичним меким гвожђем: 1 и 7. намотаји, 2. казаљка, 3. контра тег, 4, 6 и 11. јарам од меког гвожђа, 5. помично меко гвожђе, 8. држач, 9. изолатор, 10. електрични прикључак, 12. електрични спој (лем), 13. потенциометар давача}
%     \label{fig:3.99}
% \end{figure}

Електрично напајање инструмента одвија се преко прикључног завртња (10) који је причвршћен за изолациони држач (9). Струја тече кроз леви намотај (1), и у тачки 12 се грана, па преко десног намотаја (7) и потенциометра давача (13) протиче ка маси. За разлику од претходног инструмента, овде су десни намотај и потенциометар паралелно везани, чиме се смањује могућност варничења на клизачу потенциометра.

На слици 3.100 приказана је уобичајена конструкција потенциометра давача са пловком у резервоару.

% \begin{figure}[htbp]
%     \centering
%     \includegraphics[width=0.7\textwidth]{slika_3_100.png}
%     \caption{Конструкција потенциометра давача са пловком у резервоару за гориво: 1. пловак, 2. осовина, 3 и 4. доње и горње кућиште, 5. електрични прикључак, 6. спајање кућишта, 7. полуга пловка, 8. спирална опруга за довод струје, 9. клизач, 10. изолационо тело, 11. електрични извод, 12. намотаји потенциометра}
%     \label{fig:3.100}
% \end{figure}

\textbf{Електронски показивачи.} – На слици 3.101, примера ради, дата је шема везе електронског показивача, са потенциометарским давачем.
% --- Страница 208 ---
% \begin{figure}[htbp]
%     \centering
%     \includegraphics[width=0.75\textwidth]{slika_3_101.png}
%     \caption{Шема везе електронског показивача нивоа горива}
%     \label{fig:3.101}
% \end{figure}

Отпорници ($R_1$) и ($R_4$) сачињавају променљиви отпор давача у мосној вези. Излазни напон давача, појачава се помоћу симетричног појачивача са транзисторима ($T_1$) и ($T_2$). Мерни инструмент, као волтметар, везује се на уобичајен начин између два колектора. Транзистори ($T_1$) и ($T_2$), симетричног појачивача, имају заједнички отпорник у колу емитера, чиме се обезбеђује стабилан рад. Колектор транзистора ($T_3$) везује се за базу транзистора ($T_3$). У случају пораста овог колекторског напона (што означава пад нивоа горива на минимум) напон достигне пробојни напон ценер-диоде ($Z_4$). По пробоју диоде, транзистор ($T_3$) отвори и укључи црвену упозоравајућу лампицу на инструменту.

Отпорником ($R_1$) подешава се показивање инструмента при потпуно празном резервоару а отпорником ($R_2$) при потпуно пуном резервоару.

Отпорником ($R_3$) подешава се сигнализација минималног нивоа (тзв. резерве), када се пали црвена контролна лампица.

\textbf{Капацитивни давачи.} – Код појединих возила, због облика резервоара не може се поставити потенциометарски или биметални давач са пловком, јер захтевају релативно велик простор. Необичан облик резервоара често је разлог што се тражи посебно место уградње, којим се обезбеђује највећа безбедност у случају саобраћајне незгоде. У таквим случајевима користе се капацитивни или термички давачи који немају скретни пловак са полугом.

Код капацитивних давача, као мерни сигнал, користи се промена капацитета, у зависности од промене диелектрика између облога. Како је диелектрична константа горива два пута већа од ваздуха, капацитет кондензатора напуњен горивом, имаће два пута већи капацитет у односу
% --- Страница 209 ---
на онај са ваздухом. Према томе, капацитет се мења сразмерно нивоу горива у резервоару па се она може користити за мерење тог нивоа.

Сонда давача израђује се у облику тракастог кондензатора (слика 3.102а) код којих облоге не стоје једна наспрам друге, већ једна поред друге, на еластичној фолији. Повећање капацитета постиже се повећањем броја облога, које се међусобно везују у облику чешља. Посебним обликовањем тракастих облога (слика 3.102б) може се постићи линеарно показивање инструмента и при неправилном облику резервоара, у функцији нивоа и запремине.

% \begin{figure}[htbp]
%     \centering
%     \includegraphics[width=0.6\textwidth]{slika_3_102.png}
%     \caption{Капацитивни давачи: а) са равним облогама и б) са посебно обликованим облогама}
%     \label{fig:3.102}
% \end{figure}

На слици 3.103 приказана је блок-шема електронског показивача нивоа горива са капацитивним давачем. Кондензатори давача су делови осцилатора којима се одређује његова учестаност, па је излазна учестаност осцилатора сразмерна нивоу горива у резервоару. Преко фреквентног претварача овај сигнал се претвара у променљиви напон, који се преко јединица за прилагођавање одводи директно на показни инструмент. Тај исти напон, преко посебне јединице за одређивање прага вредности, може се довести и на упозоравајућу лампицу за доњи ниво горива (резерва).

Комплетна електроника показивача, са капацитивним давачем, поставља се у прикључну кутију на самом резервоару, да би се што више смањила капацитивност прикључних веза између кондензатора и електронског кола.

\textbf{Термички показивачи.} – На слици 3.104 приказана је блок-шема електронског показивача са термичким давачима у резервоару. Термички давачи су електрични отпорници који се загревају електричном струјом, чија температура зависи од тога колико су зароњени у гориво, односно од нивоа горива у резервоару. Како се њихов отпор мења са темпе-%
% --- Страница 210 ---
% \begin{figure}[htbp]
%     \centering
%     \includegraphics[width=0.75\textwidth]{slika_3_103.png}
%     \caption{Блок-шема електронског показивача нивоа горива са капацитивним давачем}
%     \label{fig:3.103}
% \end{figure}
%
% \begin{figure}[htbp]
%     \centering
%     \includegraphics[width=0.75\textwidth]{slika_3_104.png}
%     \caption{Блок-шема електронског показивача нивоа горива са термичким давачем}
%     \label{fig:3.104}
% \end{figure}
ратуром, пропуштајући кроз њих сталну струју, јасно је да ће пад напона на њиховим крајевима бити директно зависан од температуре, односно од нивоа горива. Зато се тај напон може користити за мерење нивоа. Међутим, мерени отпор зависи и од температуре горива, што се мора компензовати додавањем још једног отпорника у резервоар, који се смешта на најдебљи део резервоара да буде стално потопљен. Тако његов отпор зависи само од температуре горива а не и од нивоа. У том случају мери се практична разлика напона између ова два отпорника, па се та разлика појачивача и преко јединице за прилагођавање (у облику мерног сигнала) доводи на инструмент.

Тачност оваквог мерења, међутим, битно ремете таласања и бућкање горива у резервоару током вожње. Да би се елиминисала та негативна појава, мерна сонда се смешта у омотач у који гориво продире само кроз мали доњи отвор. На тај начин ниво горива око мерне сонде не може се брзо мењати, па таласање у резервоару не утиче на тачност мерења.

% --- Страница 211 ---
\subsection{Амперметар и волтметар на возилу}

Електрично напајање инсталације возила врши се помоћу акумулатора и алтернатора у паралелној вези. Током вожње алтернатор напаја све укључене пријемнике, и ако је потрошња мања, још допуњава и акумулатор. Када је потрошња већа од енергије коју тренутно даје алтернатор, разлика се надокнађује из акумулатора. Насупрот томе, током покретања возила и рада мотора у месту, електрично напајање се врши углавном из акумулатора. Без ова два извора, возило није у возном стању. Из тих разлога, амперметром се контролише струја пуњења, а помоћу волтметра стање акумулатора.

\subsubsection{Амперметри, шема везе, опис и конструкција}

\textbf{Шема везе.} – На возила се најчешће уграђују амперметри са нултим положајем на средини скале тако да скретање казаљке, лево или десно, показује и смер струје. Обично лево, са ознаком $(-)$, означава пражњење, а десно, са ознаком $(+)$, означава пуњење акумулатора. Такав амперметар везује се између алтернатора и акумулатора, према шеми на слици 3.105. Отпорник (R) симболично представља све електричне пријемнике на возилу, изузев стартера који се због велике струје никад не везује преко амперметра.

% \begin{figure}[htbp]
%     \centering
%     \includegraphics[width=0.5\textwidth]{slika_3_105.png}
%     \caption{Шема везе амперметра на возилу}
%     \label{fig:3.105}
% \end{figure}

\textbf{Опис и конструкција.} – На возила, због потреса и сталних померања, највише се уграђују амперметри са поларизованим меким гвожђем и дијаграмом магнетног поља.

На слици 3.106 шематски је приказана конструкција амперметра са поларизованим меким гвожђем и дијаграмом магнетног поља.

Кретни систем инструмента сачињавају: казаљка (2) и помично меко гвожђе (6), који су причвршћени на осовину (7) са лежиштима (8). Меко гвожђе се израђује од челика са мало угљеника, тако да се лако поларизује и у магнетном пољу поставља у правцу резултујућег поља. У искљученом положају, потпуно уравнотежен кретни систем, у средњем, нултом положају, држи стални магнет (4).

Када кроз инструмент протиче струја (улазни прикључак се не види на слици), око масивног проводника (види стрелице) ствара се магнетно поље. Смер линија тог поља управљан је на смер линија магнетног поља сталног магнета. Сада на, поларисано помично гвожђе, делују оба магне-%
% --- Страница 212 ---
% \begin{figure}[htbp]
%     \centering
%     \includegraphics[width=0.6\textwidth]{slika_3_106.png}
%     \caption{Конструкција амперметра са поларизованим меким гвожђем и дијаграм магнетног поља: 1. електрични прикључак, 2. казаљка, 3. магнетна облога, 4. стални магнет, 5. постоље, 6. помично меко гвожђе, 7. осовина кретног система, 8. лежиште}
%     \label{fig:3.106}
% \end{figure}
тна поља па се оно постави у правцу резултујућег поља (види слика 3.105).

Током пражњења акумулатора промени се смер струје у инструменту, па се промени и смер магнетног поља ($H_2$). То значи да поље ($H_2$) има супротан смер од поља сталног магнета ($H_1$). Зато се помично гвожђе, са казаљком, помери на супротну страну скале и амперметар показује пражњење. Угао скретања инструмента ($\varphi$), зависи од величине магнетног поља ($H_2$) коју ствара струја, па је према томе сразмеран величини струје.

Грешке услед промене температуре компензују се помоћу магнетне отоке (3), слично као код брзиномера.

\textbf{Амперметар са обртним магнетом.} – Конструкција ових инструмената слична је инструментима са укрштеним навојем и обртним магнетом. Разлика је једино у томе, што се овде, наместо малог намотаја, поставља стални магнет, који ствара магнетно поље управно у пољу главног намотаја.

На слици 3.107 приказани су принцип рада и шема везе овог инструмента.

% \begin{figure}[htbp]
%     \centering
%     \includegraphics[width=0.6\textwidth]{slika_3_107.png}
%     \caption{Принцип рада амперметра са обртним магнетом и шема везе: 1. мерни навој, 2. струјна отока (шант), 3. стални магнет, 4. обртни магнет, 5. казаљка}
%     \label{fig:3.107}
% \end{figure}

Мањи, сразмерни део мерене струје (већи део пролази кроз отоку) прелази кроз намотај инструмента (1) и ствара магнетно поље. То поље, слично описаном амперметру, сабира се са пољем сталног магнета, па обртни магнет заузима положај у смеру резултујућег поља.

Ради смањења губитака у водовима, због великих струја, отока (шант) поставља се близу акумулатора, а њен знатно мањи, али сразмеран део, води се до амперметра.
% --- Страница 213 ---
Када је инструмент без струје, на обртни магнет (4) делује само мали стални магнет (3) који га држи у средњем нултом положају. Разлика у односу на претходни инструмент је само у томе што се наместо помичног меког гвожђа налазе обртни магнети а наместо електричног вода у облику једног навоја налази се намотај (1). Помоћу отоке опсег мерења се може знатно повећати, па се могу мерити струје и до $100\text{ A}$.

\subsubsection{Волтметри, шема везе, опис и конструкција}

\textbf{Шема везе.} – Волтметри се везују паралелно са акумулатором и алтернатором и на основу тога, при искљученом мотору, види се стање акумулатора а при раду мотора показују пуњење акумулатора (слика 3.109).

Скала ових инструмената обично је подељена на жуто, зелено и црвена поља, која, посебно за искључен и укључен мотор, означавају пун или празан акумулатор, односно, недовољно, добро и превелико пуњење акумулатора (слика 3.108).

% \begin{figure}[htbp]
%     \centering
%     \includegraphics[width=0.75\textwidth]{slika_3_108.png}
%     \caption{Показивање волтметра у зависности од стања и пуњења акумулатора}
%     \label{fig:3.108}
% \end{figure}

\textbf{Опис и конструкција.} – За волтметре се највише користе инструменти са укрштеним навојем и обртним магнетом. На слици (3.110) приказана је конструкција инструмента и дијаграм магнетног поља намотаја.

У средини непомичних намотаја смештен је лако обртан магнет на истој осовини са казаљком. И код овог инструмента обртни магнет увек заузима положај у правцу резултујућег магнетног поља. Разлика је једи-%
% --- Страница 214 ---
% \begin{figure}[htbp]
%     \centering
%     \includegraphics[width=0.5\textwidth]{slika_3_109.png}
%     \caption{Шема везе волтметра у електричној инсталацији возила}
%     \label{fig:3.109}
% \end{figure}
%
% \begin{figure}[htbp]
%     \centering
%     \includegraphics[width=0.7\textwidth]{slika_3_110.png}
%     \caption{Конструкција волтметра са укрштеним навојима и обртним магнетом: 1. стални магнет, 2. обртни магнет, 3. крајеви намотаја, 4. доње тело, 5. горње тело, 6. осовина, 7. казаљка, $H_1$ – поље сталног магнета, $H_2$ – магнетно поље намотаја, $H_r$ – резултујуће магнетно поље, $\varphi$ – угао између поља сталног магнета и поља намотаја, $\alpha$ – угао пуног скретања}
%     \label{fig:3.110}
% \end{figure}
но у томе што овде поље сталног магнета у односу на поље намотаја, наместо $90^\circ$, заузима $135^\circ$ што омогућава ширу скалу од $90^\circ$ (слика 3.110).

У искљученом положају, на обртни магнет делује само поље сталног магнета ($H_1$) и он стоји у нултом положају. У укљученом положају, при највећем напону, кроз намотај протиче максимална вредност струје која ствара магнетно поље ($H_2$). Угао између тих поља је $\varphi = 135^\circ$ а угао између резултујућег поља ($H_r$) и нултог положаја је $\alpha = 90^\circ$. То је уједно и максимално скретање инструмента.

\section{Контролне лампице}

Праћење и процена вредности појединих величина, битних за рад мотора и погонску сигурност возила не поверава се пажњи возача. Тиме се смањује оптерећење возача, који више пажње може посветити условима саобраћаја и сигурности вожње, а смањује се и неопходно стручно знање за читање различитих инструмената.
% --- Страница 215 ---
Давачи, којима се прате те величине, најчешће су конструисани као прекидачи. У њима се мерене величине, најпре претварају у неке друге за поређење згодне величине, а затим пореде са датим референтним вредностима. По достизању дате вредности, укључењем контролних лампица разних боја, упозоравају возача.

Најчешће се на овај начин прати притисак уља у мотору, температура расхладне течности, пуњење акумулатора, ниво уља у кочионом систему итд. У посебну групу спадају контролне лампице које сигналишу само укључење неког уређаја.

\subsection{Контролна лампица за притисак уља у мотору}

Контролна лампица која сигналише недозвољени пад притиска уља у мотору (самим тим и недовољну количину уља) спада у најстарије контроле те врсте. Зато су и давачи за ту контролу старе конструкције који су се мало мењали све до данашњих дана.

\textbf{Давач.} – На слици (3.111) дат је пресек давача контролне лампице за притисак уља. Ради као обични уљни прекидач, који укључује и искључује своје контакте под утицајем притиска који делује на еластичну мембрану.

% \begin{figure}[htbp]
%     \centering
%     \includegraphics[width=0.7\textwidth]{slika_3_111.png}
%     \caption{Давач контролне лампице за притисак уља: 1. метално кућиште са навојем, 2. завртањ за подешавање, 3. мерна опруга, 4. изолација, 5. контра матица, 6. тањирасти контакт, 7. тањираста опруга, 8. мембрана}
%     \label{fig:3.111}
% \end{figure}

Када је притисак уља мањи од подешене вредности (обично $0{,}5\text{ bar}$), струја, преко завртња (2), опруге (3), тањирасте опруге (7) и контактне плоче (6) протиче ка кућишту (1) које се директно споји са масом. У том положају контролна лампица светли.

У случају да притисак уља порасте изнад подешене вредности, механичка сила преко мембране (8) потисне тањирасту опругу (7) и прекине
% --- Страница 216 ---
спој са контактном плочом (6) а самим тим и са масом. У том положају контролна лампица се гаси.

\subsection{Контролна лампица за температуру мотора}

Код мотора са унутрашњим сагоревањем једна од најважнијих величина која се контролише је температура мотора. То се најчешће врши преко температуре расхладне течности или уља у мотору, а код ваздушно хлађених мотора преко температуре главе мотора.

У новије време код возила са Ламбда-сондом и катализаторима, контролише се и температура издувних гасова.

Давачи су најчешће на принципу биметалних контаката, а за веће температуре на принципу промене отпорности полупроводника или промене напона термо елемената.

\textbf{Биметални давачи.} – На слици (3.112) дат је принцип рада биметалног давача са контактом који отвара и контактом који затвара при достигнутим температурама, а на слици (3.13) и конструкција таквог давача.

% \begin{figure}[htbp]
%     \centering
%     \includegraphics[width=0.6\textwidth]{slika_3_112.png}
%     \caption{Принцип рада биметалног давача са контактом који затвара и контактом који отвара при достигнутој температури}
%     \label{fig:3.112}
% \end{figure}

% \begin{figure}[htbp]
%     \centering
%     \includegraphics[width=0.6\textwidth]{slika_3_113.png}
%     \caption{Попречни пресек биметалног давача температуре (VDO): 1. изолациона чаура, 2. месингано кућиште, 3. завртањ за подешавање, 4. лепак за заливање, 5. заштита, 6. биметална трака, 7. месингана маса, 8. помоћни контакт, 9. електрични прикључак}
%     \label{fig:3.113}
% \end{figure}

% --- Страница 217 ---
Биметална трака (6) је чврсто спојена за месингано тело (7) у врху кућишта (2). Други, слободни крај траке, са помичним контактом (8), у укљученом положају, еластична сила биметала притиска, на непомични контакт (3) који уједно има и улогу завртња за подешавање. Вредност температуре укључивања подешава се овим завртњем, који се затим залива лепком да се спречи евентуално одвијање или нестручно померање. Струја се доводи преко прикључка (9) и овог завртња до непомичног контакта. Овај проводни део изолован је од кућишта изолационом чауром (1).

Овакав давач, са контактом који затвара, производи се за температуре од $92^\circ\text{C}$ до $100^\circ\text{C}$, а са контактом који отвара, за температуре од $20^\circ\text{C}$. (На слици је верзија са контактом који отвара, уцртан испрекиданим линијама). Толеранција укључивања и искључивања је максимално $5^\circ\text{C}$.

\subsection{Остале контролне лампице}

Контролне лампице се могу груписати у две групе. У прву групу спадају лампице чије упозоравајуће паљење представља најважнију информацију за возача. Паљење тих лампица понекад најављује престанак рада таквих уређаја без којих мотор не може да ради. То су већ описане лампице за притисак уља, температуру мотора, пуњење акумулатора, нестанак горива и губитак притиска у кочионом систему.

У другу групу спадају лампице које само подсећају возача да су неки уређаји укључени и раде. Ту спадају контролне лампице за дуга и оборена светла и светла за маглу, грејачи задњег стакла, ручна кочница, чок (саух) карбуратора итд.

Те лампице се углавном пале паралелно са укључењем контролисаног уређаја. Ако су то електрични уређаји, на пример светла или грејачи, струја се доводи одмах од прекидача или релеја до лампице. Ако су то механички уређаји, као на пример ручна кочница или „саух“, на механизам су уграђени мали микропрекидачи који пале контролне лампице, када се активирају ти механички уређаји.

Контролна лампица за мигавце објашњена је у поглављу о мигавцима, а контролна лампица за пуњење акумулатора и исправан рад алтернатора обрађена је у поглављу о алтернаторима.



\end{document}
